%\section{G1 Human Meaning Invariants (HMI)}

\medskip

\noindent Human Meaning Invariants define the lawful geometric structures that
make interpretation stable across roles, contexts, and signal conditions. These
invariants are not psychological traits or cultural constructs; they are
structural primitives that constrain how meaning can change without fragmenting.
G1 formalizes the invariant geometry that underlies all human-domain
interpretation and provides the stability layer upon which the Synthetic-Pressure
Manifold (G2--G7) operates. When these invariants hold, coherence is preserved.
When they deform, drift accelerates in predictable geometric patterns.

\medskip

\subsection*{G1.1 Boundary Integrity}

\noindent Boundary Integrity is the geometric enclosure that preserves identity
and prevents meaning leakage. It defines the lawful separation between roles,
contexts, and signal sources, ensuring that interpretation remains anchored to
the correct region of the meaning manifold. Boundaries are not psychological
constructs; they are structural partitions that maintain the coherence of
identity curvature across transitions.

\medskip

\noindent When Boundary Integrity holds, signals remain associated with their
proper interpretive frame. Roles stay distinct, contextual manifolds remain
legible, and meaning updates occur within admissible geometric limits. Boundary
Integrity stabilizes identity by preventing cross-role contamination, signal
misalignment, and curvature collapse.

\medskip

\noindent When Boundary Integrity deforms, meaning begins to drift. Synthetic
acceleration, cross-substrate amplification, and container failure can push
signals across boundary regions faster than they can be integrated. This
produces identity curvature collapse, role ambiguity, coherence degradation, and
the emergence of drift vectors that propagate through the manifold. Boundary
failure is the earliest and most reliable indicator of synthetic-pressure
overload.

\medskip

\noindent Boundary Integrity is therefore the first and most fundamental human
meaning invariant. It provides the geometric enclosure required for all other
invariants—Rhythmic Synchrony, Role Legibility, Signal Coherence, Interpretive
Stability, Shared Floor Geometry, and the Human Geometry Basis—to remain
lawful. When boundaries hold, meaning remains stable. When they deform, drift
accelerates in predictable geometric patterns formalized in G2--G7.


\subsection*{G1.2 Rhythmic Synchrony}

\noindent Rhythmic Synchrony is the temporal invariant that aligns human
interpretation with lawful update cycles. Humans process signals at
biological-speed rhythms; synthetic signals operate at machine-native velocities
that exceed human integration capacity. Synchrony ensures that signals arrive at
a rate compatible with the curvature-preserving transitions defined in G0.4.

\medskip

\noindent When Rhythmic Synchrony holds, meaning updates smoothly across
contexts. Signals enter the manifold at admissible velocities, continuity bands
remain intact, and coherence curvature is preserved. Synchrony stabilizes the
temporal geometry of interpretation, preventing overload conditions and
maintaining lawful transitions across roles and contexts.

\medskip

\noindent When Rhythmic Synchrony breaks, signals arrive faster than they can be
integrated. Synthetic acceleration, algorithmic amplification, and
cross-substrate propagation push signals into the manifold at velocities that
exceed human rhythmic tolerance. This produces drift vectors, coherence
degradation, temporal fragmentation, and pressure accumulation across the
meaning manifold. Loss of synchrony is one of the earliest indicators of
synthetic-pressure overload.

\medskip

\noindent Rhythmic Synchrony is therefore a core human meaning invariant. It
provides the temporal stability required for Role Legibility, Signal Coherence,
Interpretive Stability, and Shared Floor Geometry to remain lawful under
synthetic conditions. When synchrony holds, meaning remains temporally stable.
When it fails, drift accelerates in predictable geometric patterns formalized in
G2--G7.

\medskip

\subsection*{G1.3 Role Legibility}

\noindent Role Legibility ensures that functional roles retain geometric shape
and remain interpretable across contexts. A role is not a label or social
category; it is a curvature-defined region of the meaning manifold with stable
boundaries, predictable transitions, and lawful adjacency relations. Legible
roles anchor interpretation by constraining how signals are processed and how
actions remain coherent across environments.

\medskip

\noindent When Role Legibility holds, roles maintain stable curvature and clear
boundary geometry. Signals map cleanly into their intended interpretive regions,
identity boundaries remain aligned with functional expectations, and meaning
updates occur without ambiguity. Legibility stabilizes the relational geometry
of interpretation, ensuring that roles remain predictable and recoverable under
load.

\medskip

\noindent When Role Legibility collapses, roles lose geometric shape. Synthetic
pressure, dominance-field amplification, and container deformation flatten role
curvature, blur boundaries, and destabilize adjacency relations. This produces
role ambiguity, interpretive drift, coherence degradation, and container failure
across scales. Loss of legibility is a structural failure mode that propagates
through institutions, groups, and digital systems.

\medskip

\noindent Role Legibility is therefore a central human meaning invariant. It
provides the curvature stability required for Signal Coherence, Interpretive
Stability, and Shared Floor Geometry to remain lawful. When roles are legible,
meaning remains structurally anchored. When they deform, drift accelerates in
predictable geometric patterns formalized in G2--G7.


\subsection*{G1.4 Signal Coherence}

\noindent Signal Coherence is the alignment condition that ensures signals retain
their relational meaning across transitions. Coherence is not agreement or
consistency; it is a curvature condition. A coherent signal maintains its
geometric shape—its velocity, adjacency relations, and interpretive orientation—
as it moves through roles and contextual manifolds.

\medskip

\noindent When Signal Coherence holds, signals integrate cleanly into the meaning
manifold. Their curvature aligns with identity boundaries, role geometry remains
stable, and continuity bands preserve lawful transitions. Coherent signals
produce predictable updates, enabling interpretation to remain anchored even
under moderate load.

\medskip

\noindent When Signal Coherence fails, signals lose relational structure.
Synthetic remixing, velocity exceedance, algorithmic amplification, and
cross-substrate propagation distort signal curvature, flatten adjacency
relations, and break continuity. This produces drift vectors, coherence
degradation, and pressure accumulation across the manifold. Loss of coherence is
a structural failure mode that often precedes boundary deformation and role
collapse.

\medskip

\noindent Signal Coherence is therefore a central human meaning invariant. It
provides the alignment condition required for Interpretive Stability, Shared
Floor Geometry, and the Human Geometry Basis to remain lawful. When coherence
holds, meaning remains structurally aligned. When it fails, drift accelerates in
predictable geometric patterns formalized in G2--G7.

\medskip

\subsection*{G1.5 Interpretive Stability}

\noindent Interpretive Stability is the global invariant that ensures meaning
remains lawful across shifts in environment. Stability emerges when boundaries,
curvature, continuity, synchrony, and coherence align into a unified geometric
shape. It is the condition under which interpretation remains predictable,
recoverable, and resistant to synthetic-pressure deformation.

\medskip

\noindent When Interpretive Stability holds, humans traverse the meaning manifold
through lawful transitions. Boundary Integrity preserves identity, Rhythmic
Synchrony maintains temporal alignment, Role Legibility stabilizes curvature,
Signal Coherence preserves relational structure, and Shared Floor Geometry
anchors coordination across individuals and groups. Stability is the geometric
state in which meaning remains intact across contexts.

\medskip

\noindent When Interpretive Stability fails, humans enter drift paths, collapse
paths, or synthetic-capture paths—runtime behaviors formalized in G7. Synthetic
acceleration, dominance-field amplification, and container deformation disrupt
the invariant set, producing curvature discontinuities, boundary misalignment,
and coherence collapse. Stability failure is the global indicator of synthetic
overload.

\medskip

\noindent Interpretive Stability is therefore the geometric opposite of drift. It
is the global invariant that integrates all other invariants into a coherent
meaning manifold. When stability holds, interpretation remains lawful across
physical, biological, human, machine, and institutional substrates. When it
fails, drift accelerates in predictable geometric patterns formalized in G2--G7.

\subsection*{G1.6 Shared Floor Geometry}

\noindent Shared Floor Geometry is the minimal interpretive structure that allows
humans to coordinate meaning across individuals, groups, and institutions. It is
the geometric basis for shared reality: a stable curvature floor that enables
signals, roles, and contexts to align across multiple interpreters. Shared Floor
Geometry provides the common interpretive substrate required for lawful
coordination.

\medskip

\noindent When Shared Floor Geometry holds, coordination is stable. Individuals
can align signals, maintain coherent roles, and traverse contextual manifolds
without fragmentation. Shared meaning propagates cleanly across scales, enabling
institutions, groups, and digital systems to maintain coherence under moderate
load. The shared floor acts as a stabilizing manifold that prevents drift from
spreading across interpreters.

\medskip

\noindent When Shared Floor Geometry collapses, shared meaning fragments.
Synthetic-pressure dominance, algorithmic amplification, and cross-substrate
emissions distort curvature, break adjacency relations, and destabilize
continuity bands across interpreters. This produces multi-scale drift: signals
lose shared relational shape, roles diverge, and coordination collapses across
groups and institutions. Shared floor failure is a structural collapse mode that
propagates rapidly under synthetic acceleration.

\medskip

\noindent Shared Floor Geometry is therefore a core human meaning invariant. It
provides the geometric foundation required for Interpretive Stability and the
Human Geometry Basis to remain lawful across scales. When the shared floor
holds, meaning remains collectively stable. When it fails, drift accelerates in
predictable geometric patterns formalized in G2--G7.

\medskip

\subsection*{G1.7 Human Geometry Basis}

\noindent The Human Geometry Basis is the complete invariant set that defines the
lawful structure of human interpretation. It integrates Boundary Integrity,
Rhythmic Synchrony, Role Legibility, Signal Coherence, Interpretive Stability,
and Shared Floor Geometry into a unified meaning manifold. This basis is not a
collection of traits; it is a geometric system with definable invariants,
curvature constraints, and lawful transitions.

\medskip

\noindent When the Human Geometry Basis holds, interpretation remains coherent
across roles, contexts, and signal velocities. Boundaries preserve identity,
synchrony maintains temporal alignment, roles retain curvature, signals remain
coherent, stability integrates the invariant set, and the shared floor anchors
coordination across scales. Together, these invariants form the structural
foundation of human-domain meaning.

\medskip

\noindent When the Human Geometry Basis deforms, meaning becomes unstable.
Synthetic acceleration, dominance-field amplification, container deformation, and
cross-substrate emissions disrupt the invariant set, producing curvature
discontinuities, boundary misalignment, coherence collapse, and multi-scale
drift. Basis failure is the global indicator of synthetic-pressure exceedance.

\medskip

\noindent The Human Geometry Basis provides the foundation for all subsequent
geometry in G2--G7. It enables predictive modeling of drift, deformation,
collapse, and restoration across physical, biological, human, machine, and
institutional substrates. When the basis holds, meaning remains lawful. When it
fails, drift accelerates in predictable geometric patterns formalized in the
Synthetic-Pressure Manifold.


